\documentclass[11pt,a4paper]{article}
\usepackage[utf8]{inputenc}
\usepackage[T1]{fontenc}
\usepackage{geometry}
\geometry{margin=1in}
\usepackage{hyperref}
\usepackage{enumitem}
\usepackage{booktabs}
\usepackage{longtable}
\usepackage{xcolor}
\usepackage{titlesec}
\usepackage{amsmath,amssymb}

\titleformat{\section}{\Large\bfseries\color{blue!70!black}}{}{0em}{}
\titleformat{\subsection}{\large\bfseries\color{blue!50!black}}{}{0em}{}
\titleformat{\subsubsection}{\normalsize\bfseries\color{blue!30!black}}{}{0em}{}

\title{\textbf{Replication Plan}\\[0.5em]
\large Simple Coplanar Waveguide Resonator Mask Targeting Metal--Substrate Interface\\[0.3em]
\normalsize OSTI ID: 1983793}
\author{Auto-generated Replication Blueprint}
\date{\today}

\begin{document}
\maketitle
\tableofcontents
\newpage

%---------------------------------------------------------------------
\section{Paper Overview}
%---------------------------------------------------------------------
This paper presents the design and finite-element electromagnetic simulation of coplanar waveguide (CPW) resonators for superconducting quantum computing, with mask geometries specifically targeting the metal--substrate (MS) interface to study surface dielectric losses. The authors vary geometric parameters (center conductor width, gap width, film thickness) and compute electric field energy participation ratios in different loss regions (MS interface, substrate--air interface, metal--air interface, bulk substrate). The quality factors and resonance frequencies are then predicted from these participation ratios and compared with experimental measurements.

%---------------------------------------------------------------------
\section{Detailed Workflow}
%---------------------------------------------------------------------

\subsection{Phase 1: Geometry Definition}
\begin{enumerate}[label=\textbf{1.\arabic*}]
  \item \textbf{Define CPW resonator geometry.}
    \begin{itemize}
      \item Center conductor width $w$, gap $s$, film thickness $t$, substrate thickness.
      \item Parametric sweep values as specified in the paper.
      \item Ground plane extent, coupling structures (if modeled).
    \end{itemize}
  \item \textbf{Define interface/surface loss layers.}
    \begin{itemize}
      \item Thin dielectric layers at MS, SA, MA interfaces (thickness $\sim 3$\,nm).
      \item Assign dielectric constants and loss tangents for each layer.
    \end{itemize}
\end{enumerate}

\subsection{Phase 2: Finite-Element Simulation}
\begin{enumerate}[label=\textbf{2.\arabic*}]
  \item \textbf{Set up the 2D cross-section model} in a FEM solver.
    \begin{itemize}
      \item Eigenmode analysis to find the CPW quasi-TEM mode.
      \item Adaptive mesh refinement near interfaces.
    \end{itemize}
  \item \textbf{Compute electric field distribution} $\mathbf{E}(x, y)$.
  \item \textbf{Compute participation ratios.}
    \begin{equation*}
      p_i = \frac{\int_{\text{region}_i} \varepsilon_i |\mathbf{E}|^2 \, dA}{\int_{\text{total}} \varepsilon |\mathbf{E}|^2 \, dA}
    \end{equation*}
  \item \textbf{Predict quality factor.}
    \begin{equation*}
      \frac{1}{Q} = \sum_i p_i \tan \delta_i
    \end{equation*}
\end{enumerate}

\subsection{Phase 3: Parametric Sweep}
\begin{enumerate}[label=\textbf{3.\arabic*}]
  \item Sweep $w$ and $s$ over specified ranges.
  \item For each geometry, compute participation ratios and predicted $Q$.
  \item \textbf{Reproduce key plots}: $p_{\text{MS}}$ vs.\ $w/s$ ratio, $Q$ vs.\ geometry.
\end{enumerate}

\subsection{Phase 4: Comparison with Experiment}
\begin{enumerate}[label=\textbf{4.\arabic*}]
  \item Compare simulated $Q$ values and resonance frequencies with experimental data from the paper.
  \item Extract effective loss tangents from fitting.
\end{enumerate}

%---------------------------------------------------------------------
\section{Datasets}
%---------------------------------------------------------------------

\begin{longtable}{p{4.5cm} p{3.5cm} p{6cm}}
\toprule
\textbf{Dataset / Source} & \textbf{Type} & \textbf{Location / Notes} \\
\midrule
\endfirsthead
\toprule
\textbf{Dataset / Source} & \textbf{Type} & \textbf{Location / Notes} \\
\midrule
\endhead
CPW geometry parameters & From paper & Tables of $w$, $s$, $t$ values \\
\midrule
Material properties & Literature & Dielectric constants and loss tangents for Si, SiO$_2$, Al, Nb \\
\midrule
Experimental $Q$ / $f_0$ data & From paper & Tables/figures of measured resonator properties \\
\bottomrule
\end{longtable}

%---------------------------------------------------------------------
\section{Models, Tools, and Software}
%---------------------------------------------------------------------

\begin{longtable}{p{4.5cm} p{2.5cm} p{6.5cm}}
\toprule
\textbf{Tool / Library} & \textbf{Version} & \textbf{Purpose / URL} \\
\midrule
\endfirsthead
\toprule
\textbf{Tool / Library} & \textbf{Version} & \textbf{Purpose / URL} \\
\midrule
\endhead
COMSOL Multiphysics & 6.0+ & FEM electromagnetic eigenmode simulation \newline \url{https://www.comsol.com/} \\
\midrule
Ansys HFSS & latest & Alternative FEM EM solver \newline \url{https://www.ansys.com/products/electronics/ansys-hfss} \\
\midrule
Sonnet & latest & Alternative planar EM simulator \newline \url{https://www.sonnetsoftware.com/} \\
\midrule
PALACE (AWS) & latest & Open-source FEM for superconducting circuits \newline \url{https://github.com/awslabs/palace} \\
\midrule
Gmsh & $\geq 4.10$ & Open-source mesh generation \newline \url{https://gmsh.info/} \\
\midrule
FEniCS / FEniCSx & latest & Open-source FEM framework \newline \url{https://fenicsproject.org/} \\
\midrule
Python (NumPy, Matplotlib) & $\geq 3.8$ & Post-processing and plotting \\
\midrule
Qiskit Metal & latest & Optional: superconducting qubit design \newline \url{https://qiskit.org/metal/} \\
\bottomrule
\end{longtable}

\subsection{Hardware Requirements}
\begin{itemize}
  \item 2D cross-section FEM: any workstation; $< 1$\,min per geometry.
  \item 3D full-resonator simulations: $\geq$ 32\,GB RAM, multi-core.
\end{itemize}

\end{document}
